\documentclass[10pt,a3paper]{article}

\usepackage[a3paper, left=20mm, right=60mm, top=15mm, bottom=15mm]{geometry}
\usepackage{fontspec}
\setmainfont{TeX Gyre Heros}
\usepackage{amsmath,amssymb,mathtools}
\usepackage{array}

\setlength{\parskip}{2pt}
\setlength{\abovedisplayskip}{4pt}
\setlength{\belowdisplayskip}{4pt}
\setlength{\abovedisplayshortskip}{2pt}
\setlength{\belowdisplayshortskip}{2pt}
\setlength{\extrarowheight}{2pt}

\pagestyle{empty}

% top-aligned ragged-right p column
\newcolumntype{T}[1]{>{\raggedright\arraybackslash}p{#1}}

\usepackage{eso-pic}
\usepackage{tikz}
\usetikzlibrary{calc}
\usepackage[useregional]{datetime2}

% Running margin label
\newcommand{\mymarginlabel}{\textcopyright\ David Ewing dew59@uclive.ac.nz | \DTMnow\ | \jobname.tex}
\AddToShipoutPictureBG{%
  \begin{tikzpicture}[remember picture,overlay]
    \node[rotate=90, anchor=north]
      at ($(current page.north west)+(1.4cm,-13cm)$)
      {\scriptsize\ttfamily \mymarginlabel};
  \end{tikzpicture}%
}

\begin{document}

\noindent
\begin{tabular}{@{}T{0.46\linewidth}|T{0.46\linewidth}@{}}
\hline
\textbf{Newton on KKT} & \textbf{KKT Analytical} \\
\hline\hline
% row 1 — problem (both)
$\min\ f(x)=x_1^2+5x_2^2$,\; s.t.\ $g(x)=4-x_1-x_2\leq 0$
&
$\min\ f(x)=x_1^2+5x_2^2$,\; s.t.\ $g(x)=x_1+x_2-4\leq 0$
\\[3pt]\hline
% row 2 — slack (KKT only)
&
\textbf{Slack:} $g(x)-s^2=0$,\; so $l(x,\lambda,s)=f(x)-\lambda(g(x)-s^2)$
\\[3pt]\hline
% row 3 — dl/ds (KKT only)
&
$\dfrac{\partial l}{\partial s}=2\lambda s=0 \;\Rightarrow\; \lambda=0$ (inactive) or $s=0$ (active)
\\[3pt]\hline
% row 4 — test (KKT only)
&
\textbf{Test:} $X_\text{unc}=(0,0)$;\; $g(X_\text{unc})=-4<0$ \quad $\therefore$ infeasible $\Rightarrow$ $\lambda\neq 0$
\quad $\xRightarrow{2\lambda s=0}$ \quad $s=0$, constraint \textbf{active}
\\[3pt]\hline
% row 5 — Lagrangian (both)
$l(x,\lambda)=x_1^2+5x_2^2-\lambda(x_1+x_2-4)$ \quad (same $l$, implicit in $\nabla_x l$)
&
$l(x,\lambda)=x_1^2+5x_2^2-\lambda(x_1+x_2-4)$ \quad (set $s=0$)
\\[3pt]\hline
% row 6 — r and J (both, first parallel)
\begin{minipage}[t]{\linewidth}
Assert $r(x,\lambda)=0$ and $J$; note $J=\begin{bmatrix}H_l & -\nabla_x g \\ \nabla_x g^\top & 0\end{bmatrix}$:
\[
  r = \begin{bmatrix} \nabla_x l \\ g(x) \end{bmatrix}
    = \begin{bmatrix} 2x_1-\lambda \\ 10x_2-\lambda \\ x_1+x_2-4 \end{bmatrix}
  \quad
  J = \begin{bmatrix} 2&0&-1\\0&10&-1\\1&1&0 \end{bmatrix}
\]
\end{minipage}
&
\begin{minipage}[t]{\linewidth}
Same $r$ and $J$ arrived at from $\nabla_x l=0$ and $g=0$:
\[
  r = \begin{bmatrix} \nabla_x l \\ g(x) \end{bmatrix}
    = \begin{bmatrix} 2x_1-\lambda \\ 10x_2-\lambda \\ x_1+x_2-4 \end{bmatrix}
  \quad
  J = \begin{bmatrix} 2&0&-1\\0&10&-1\\1&1&0 \end{bmatrix}
\]
\end{minipage}
\\\hline
% row 7 — init (Newton only)
\begin{minipage}[t]{\linewidth}
\textbf{Init:} $z_0=(0,0,0)^\top$; evaluate $r$ at $z_0$:
\[
  r(z_0) = \begin{bmatrix}2x_1-\lambda\\10x_2-\lambda\\x_1+x_2-4\end{bmatrix}_{z_0=(0,0,0)}
  = \begin{bmatrix}2(0)-0\\10(0)-0\\0+0-4\end{bmatrix}
  = \begin{bmatrix}0\\0\\-4\end{bmatrix}
\]
\end{minipage}
&
{\footnotesize (no iterate — constraint known active from test above)}
\\[3pt]\hline
% row 8 — stationarity conditions (both, parallel)
\begin{minipage}[t]{\linewidth}
\textbf{Newton step:} $J\,\Delta z=-r(z_0)$, i.e.\ $-r(z_0)=(0,0,4)^\top$:
\[
  J\,\Delta z =
  \begin{bmatrix}2&0&-1\\0&10&-1\\1&1&0\end{bmatrix}
  \begin{bmatrix}\Delta x_1\\\Delta x_2\\\Delta\lambda\end{bmatrix}
  =\begin{bmatrix}0\\0\\4\end{bmatrix}
\]
Row 1: $2\Delta x_1-\Delta\lambda=0 \;\Rightarrow\; \Delta x_1=\tfrac{\Delta\lambda}{2}$

Row 2: $10\Delta x_2-\Delta\lambda=0 \;\Rightarrow\; \Delta x_2=\tfrac{\Delta\lambda}{10}$
\end{minipage}
&
\begin{minipage}[t]{\linewidth}
$\nabla_x l = H_f\,x - \lambda\,\nabla_x g = 0 \;\Rightarrow\; x^*(\lambda)=H_f^{-1}\lambda\,\nabla_x g$:
\[
  \begin{bmatrix} 2x_1-\lambda \\ 10x_2-\lambda \end{bmatrix} = 0
  \;\Rightarrow\;
  x^*(\lambda) = \begin{bmatrix} \lambda/2 \\ \lambda/10 \end{bmatrix}
\]
\end{minipage}
\\\hline
% row 9 — constraint substitution (both, parallel)
Row 3: $\tfrac{\Delta\lambda}{2}+\tfrac{\Delta\lambda}{10}=4
  \;\Rightarrow\; 6\Delta\lambda=40
  \;\Rightarrow\; \Delta\lambda=\tfrac{20}{3}$
&
Sub $x^*(\lambda)$ into $g=0$:\;
$\tfrac{\lambda}{2}+\tfrac{\lambda}{10}=4
  \;\Rightarrow\; 10\lambda+2\lambda=80
  \;\Rightarrow\; \lambda=\tfrac{20}{3}$
\\[3pt]\hline
% row 10 — solution values (both, parallel)
$\Delta x_1=\dfrac{10}{3}$,\quad $\Delta x_2=\dfrac{2}{3}$;\quad
$z_1=z_0+\Delta z=\!\left(\tfrac{10}{3},\,\tfrac{2}{3},\,\tfrac{20}{3}\right)$
&
$x^*=\!\left(\dfrac{10}{3},\,\dfrac{2}{3}\right)^\top$
\\[3pt]\hline
% row 11 — objective (both, parallel)
$f(x^*)=\!\left(\tfrac{10}{3}\right)^2+5\!\left(\tfrac{2}{3}\right)^2
  =\tfrac{100}{9}+\tfrac{20}{9}=\tfrac{40}{3}$
&
$f(x^*)=\!\left(\tfrac{10}{3}\right)^2+5\!\left(\tfrac{2}{3}\right)^2
  =\tfrac{100}{9}+\tfrac{20}{9}=\tfrac{40}{3}$
\\\hline
\end{tabular}

\vspace{8pt}

\begin{minipage}[t]{0.46\linewidth}

\noindent\textbf{SQP connection}

\noindent\textbf{KKT} — necessary conditions for a constrained optimum; gives the system $r(x,\lambda)=0$

\noindent\textbf{Lagrangian} — function whose stationarity encodes the KKT conditions; introduces $\lambda$

\noindent\textbf{Slack variables} — convert inequalities to equalities; lead to complementary slackness $\lambda s = 0$

\noindent\textbf{SQP} — solves $r=0$ using Newton's method; each step inverts the Jacobian of $r$:
\[
  r(x,\lambda) = \begin{bmatrix} \nabla_x \mathcal{L} \\ g(x) \end{bmatrix}
               = \begin{bmatrix} 2x_1-\lambda \\ 10x_2-\lambda \\ x_1+x_2-4 \end{bmatrix} = 0
  \qquad
  J = \frac{\partial r}{\partial(x_1,x_2,\lambda)}
    = \begin{bmatrix} 2 & 0 & -1 \\ 0 & 10 & -1 \\ 1 & 1 & 0 \end{bmatrix}
\]

\noindent\textbf{BFGS} — quasi-Newton approximation of the Hessian of $\mathcal{L}$ when exact computation is too costly

\end{minipage}

\vspace{8pt}

\noindent\rule{\linewidth}{0.4pt}

\vspace{4pt}

\noindent\textbf{Newton's method applied to KKT} (same example: $\min\ x_1^2+5x_2^2$ s.t.\ $x_1+x_2=4$, starting at $z_0=(0,0,0)$)

\vspace{4pt}

\noindent
\begin{minipage}[t]{0.30\linewidth}

\textbf{Step 1 — residual at $z_0$:}
\[
  r(z_0) = \begin{bmatrix} 2(0)-0 \\ 10(0)-0 \\ 0+0-4 \end{bmatrix} = \begin{bmatrix} 0 \\ 0 \\ -4 \end{bmatrix}
\]

\textbf{Step 2 — Jacobian} (constant here since $r$ is linear):
\[
  J = \begin{bmatrix} 2 & 0 & -1 \\ 0 & 10 & -1 \\ 1 & 1 & 0 \end{bmatrix}
\]

\end{minipage}\hfill
\begin{minipage}[t]{0.64\linewidth}

\textbf{Step 3 — solve} $J\,\Delta z = -r(z_0)$:
\[
  \begin{bmatrix} 2 & 0 & -1 \\ 0 & 10 & -1 \\ 1 & 1 & 0 \end{bmatrix}
  \begin{bmatrix} \Delta x_1 \\ \Delta x_2 \\ \Delta\lambda \end{bmatrix}
  = \begin{bmatrix} 0 \\ 0 \\ 4 \end{bmatrix}
\]

Row 1: $2\Delta x_1 - \Delta\lambda = 0 \;\Rightarrow\; \Delta x_1 = \tfrac{\Delta\lambda}{2}$

Row 2: $10\Delta x_2 - \Delta\lambda = 0 \;\Rightarrow\; \Delta x_2 = \tfrac{\Delta\lambda}{10}$

Row 3: $\Delta x_1 + \Delta x_2 = 4 \;\Rightarrow\; \tfrac{\Delta\lambda}{2} + \tfrac{\Delta\lambda}{10} = 4 \;\Rightarrow\; 6\Delta\lambda = 40 \;\Rightarrow\; \Delta\lambda = \tfrac{20}{3}$

\[
  \Delta x_1 = \frac{10}{3} \qquad \Delta x_2 = \frac{2}{3}
\]

\textbf{Step 4 — update:} $z_1 = z_0 + \Delta z$
\[
  z_1 = \left(\frac{10}{3},\ \frac{2}{3},\ \frac{20}{3}\right) = x^* \quad \checkmark \quad \text{(one step: KKT system is linear)}
\]

\end{minipage}

\end{document}
